\documentclass[10pt,a4paper]{amsart}
\usepackage[margin=19mm]{geometry}
\usepackage{amsmath,amssymb,mathtools}
\usepackage[scr=rsfs,cal=euler]{mathalfa}
\usepackage{xfrac}
\usepackage[unicode,hidelinks,bookmarks=false]{hyperref}
\newcommand{\ie}{i.e.\ }
\newcommand{\cf}{cf.\ }
\newcommand{\resp}{respectively\ }
\newcommand{\Subsection}{\subsection}
\providecommand{\nobreakdash}{\nobreak\hbox{-}}
\setlength{\emergencystretch}{2em}
\multlinegap0pt
\usepackage{csquotes}
\usepackage{amssymb}
\usepackage{xcolor}
\newtheorem{definition}{Definition}
\newtheorem{prop}{Proposition}
\newtheorem{lemma}{Lemma}
\newtheorem{theorem}{Theorem}
\def\Ind{\setbox0=\hbox{$x$}\kern\wd0\hbox to 0pt{\hss$\mid$\hss} \lower.9\ht0\hbox to 0pt{\hss$\smile$\hss}\kern\wd0}
\def\Notind{\setbox0=\hbox{$x$}\kern\wd0\hbox to 0pt{\mathchardef \nn=12854\hss$\nn$\kern1.4\wd0\hss}\hbox to 0pt{\hss$\mid$\hss}\lower.9\ht0 \hbox to 0pt{\hss$\smile$\hss}\kern\wd0}
\def\ind{\mathop{\mathpalette\Ind{}}}
\def\nind{\mathop{\mathpalette\Notind{}}}
\title{Simple Homogeneous Structures and Indiscernible Sequence Invariants}
\author{John Baldwin, James Freitag, Scott Mutchnik}
\date{Source: arXiv:2405.08211v3 (CC BY 4.0)}
\begin{document}
\maketitle

\noindent\emph{Source and licence.} Simple Homogeneous Structures and Indiscernible Sequence Invariants. Version arXiv:2405.08211v3 (CC BY 4.0), distributed by arXiv under the Creative Commons Attribution 4.0 International licence, http://creativecommons.org/licenses/by/4.0/, as recorded in the arXiv licence metadata for this exact version and shown on the abstract page https://arxiv.org/abs/2405.08211v3. Full licence notice: \texttt{licenses/arxiv-2405.08211-CC-BY-4.0.txt}.\par\smallskip

\noindent\textit{Editorial scope.} Counted unit: the article's theorem theorem (source0.tex lines 2941--2943). The article's complete original proof of it is delivered with it (source0.tex lines 2945--2984, one \texttt{proof} environment of 40 lines). No mathematics is written, repaired or re-derived: the statement and the proof are the article's own lines, in the article's own notation, and no retained line is substituted or re-flowed.  Retained uncounted as the source-local context the proof needs: lines 672--680; lines 688--708; lines 2486--2488.  Nothing was omitted from any retained span, so the package carries no omission record.  No retained line contains a citation, so the wrapper carries no bibliography and no bibliography entry.  No retained line contains a figure environment, an image-inclusion command or drawing code, so no image is delivered and the package declares no asset dependency.  Editorial formatting only, in addition: an AMS \texttt{amsart} preamble that restates the article's own declarations and macros used by the retained lines, namely the environments \texttt{definition}, \texttt{prop}, \texttt{lemma}, \texttt{theorem} on separate counters, together with the standard packages those lines need.  The retained environments print in this document as Definition 1, Proposition 1, Lemma 1, Theorem 1 in order of appearance, whereas the article prints the same items with the numbering of its own full document.  The complete native source of the article as distributed in the arXiv e-print for this exact version is bundled unchanged as \texttt{sources/159099/source0.tex}.
\begin{definition}
   \label{resolvable}
    
    (1) Let $I$ be an indiscernible sequence over $\emptyset$. Then $I$ is \emph{$\kappa$-resolvable} if there is some set $A \subset \mathbb{M}$ (i.e. of real elements), $|A| \leq \kappa$ so that $I$ is a Morley sequence over $A$.
   
    (2) Let $p \in S(B)$ be a type. Then $p$ is \emph{$\kappa$-resolvable} if there is some set $A \subset \mathbb{M}$, $|A| \leq \kappa$ and $ q \in S(AB)$, $q \supseteq p$ so that $q$ forks over neither $A$ nor $B$.

    
\end{definition}

\begin{prop}\label{resolvabilitygrank}
    Let $T$ be simple, and let $I$ be an indiscernible sequence. Then the following are equivalent:

    (1) The indiscernible sequence $I$ is $\kappa$-resolvable.

    (2) For some type $p \in S(B)$, so that $I$ is a Morley sequence in $p$ over $B$, $p$ is $\kappa$-resolvable.

    (2') For every type $p \in S(B)$, so that $I$ is a Morley sequence in $p$ over $B$, $p$ is $\kappa$-resolvable.

    (2'') The type $\mathrm{lim}^{+}(I/I) \in S(I)$ is $\kappa$-resolvable.

    (3) For some type $p \in S(B)$, so that $I$ is a Morley sequence in $p$ over $B$, $G(\mathrm{Cb}(p)) \leq n$.


    (3') For every type $p \in S(B)$, so that $I$ is a Morley sequence in $p$ over $B$, $G(\mathrm{Cb}(p)) \leq n$.


    (3'') $G(\mathrm{Cb}(\mathrm{lim}^{+}(I/I))) \leq n$

    
\end{prop}

\begin{lemma}\label{forkingchains}
    Let $T$ be simple, but not supersimple. Then there is a limit ordinal $\lambda$, a finite tuple $a$, and sets $B_{\gamma}$, $\gamma < \lambda$, with $B_{\gamma} \subset B_{\gamma'} $ and $a \nind_{B_{\gamma}} B_{\gamma'}$ for $\gamma < \gamma' < \lambda$, so that for $B = \bigcup_{\lambda < \gamma} B_{\gamma}$, $\mathrm{SU}(\mathrm{tp}(a/B)) \in \mathrm{Ord}$.
\end{lemma}

\begin{theorem}
    Let $T$ be simple, and let $\mathrm{SU}(\mathrm{tp}(b/\emptyset))=\infty$. Then there are ($\emptyset$-)indiscernible sequences $I^{i}=(b^{i}_{j})_{j < \omega}$, for $i < \omega$, with $b^{i}_{0} = b$, so that $I^{i} \not\equiv I^{j}$ for $i \neq j < \omega$.
\end{theorem}

\begin{proof}
    As in Definition \ref{resolvable}, we define the following invariants of indiscernible sequences.

    \begin{definition}
        \label{qresolvable}

        (1) Let $r(x) \in S(C)$ be a complete type, and let $I=\{a_{i}\}_{i < \omega}$ be an indiscernible sequence over $\emptyset$ so that $|a_{i}|= x$. Then $I$ is \textit{r}-resolvable if there is some automorphism $\sigma$ of $\mathbb{M}$, so that for $A = \sigma(C)$, $I$ is a Morley sequence over $A$ of realizations of $\sigma(r)$.  

        (2) Let $p(x) \in S(B)$ be a type. Then $p$ is \emph{$r$-resolvable} if there is some automorphism $\sigma$ of $\mathbb{M}$ so that, for $B = \sigma(C)$, there is some  $ q \in S(AB)$, $q \supseteq p, \sigma(r)$ so that $q$ forks over neither $A$ nor $B$.
    \end{definition}


Then similarly to Proposition \ref{resolvabilitygrank}, the following holds:

\begin{prop}\label{resolvabilitygrank2}
    Let $T$ be simple, and let $I$ be an indiscernible sequence. Then the following are equivalent:

    (1) The indiscernible sequence $I$ is $r$-resolvable.

    (2) For some type $p \in S(B)$, so that $I$ is a Morley sequence in $p$ over $B$, $p$ is $r$-resolvable.

    (2') For every type $p \in S(B)$, so that $I$ is a Morley sequence in $p$ over $B$, $p$ is $r$-resolvable.

    (2'') The type $\mathrm{lim}^{+}(I/I) \in S(I)$ is $r$-resolvable.

    
\end{prop}

Let $p_{0}(x) = \mathrm{tp}(b)$. It suffices to show that there do not exist finitely many $r_{1}, \ldots, r_{n}$ so that every extension of $p_{0}(x)$ is $r_{i}$-resolvable for some $1 \leq i \leq n$: otherwise, since a type $p(x)$ is always $p$-resolvable, we can find types $p_{i}(x) \in S(A_{i})$ extending $p_{0}$, for $i < \omega$, so that for $i \neq j$, $p_{i}(x)$ is not resolvable with respect to all of the same types as $p_{j}(x)$. By the previous proposition, for $I_{i}$ a Morley sequence over $A_{i}$ consisting of realizations of $p_{i}$, $I_{i}$ is resolvable with respect to exactly the same types as $p_{i}$, so for $i \neq j$, $I_{i} \not \equiv I_{j}$ because resolvability with respect to any given type is an invariant property of an indiscernible sequence.


Suppose otherwise, so such $r_{i} \in S(C_{i})$ do exist for $1 \leq i \leq n$. Because $\mathrm{tp}(b) $ has $\mathrm{SU}$-rank $\infty$, we may find $a$, $B$ as Lemma \ref{forkingchains}, which does not assume $\omega$-categoricity, and additionally so that $\mathrm{tp}(a/\emptyset)=p_{0}$. So $\mathrm{tp}(a/B)$ is ranked but for any finite $B_{0} \subset B$, $\mathrm{SU}(\mathrm{tp}(a/B_{0}))= \infty$. Without loss of generality, we assume that for some $1 \leq m \leq n$ and all $1 \leq i \leq m$, there is some finite $B_{0} \subset B$ so that $\mathrm{tp}(a/B_{0})$ is $r_{i}$-resolvable, and moreover that for any finite $B_{0} \subset B$, $\mathrm{tp}(a/B_{0})$ is $r_{i}$-resolvable for some $1 \leq i \leq m$.   So for $1 \leq i \leq m$, $\mathrm{SU}(r_{i})= \infty$.

Let $Y$ be a set of variables large enough to enumerate all of the $C_{i}$. For $c \models r_{i} \in S(C_{i})$, let $s_{i}(Y, x)= \mathrm{tp}( C_{i}, c/\emptyset)$; for $|Z| = |B|$, let $\tilde{\Phi}_{i}(Z, C_{i}, c) $ denote the set of formulas $\varphi(Z, C_{i}, c)$ that do not fork over  $C_{i}$. Let $\Phi_{i}(Z, Y, x)$ denote the set of conjunctions of formulas from $s_{i}(Y, x) \cup \tilde{\Phi}(Z, Y, x)$.  We show that the set

$$\{\vee^{m}_{i=1} \varphi_{i}(B, Y, a): \varphi_{i} \in \Phi_{i}(B, Y, a)\}$$

is consistent. This will be sufficient, because by compactness, one of the $\Phi_{i}(B, Y ,a)$ will then be consistent, which means that there will be some automorphism $\sigma$ of $\mathbb{M}$ so that, for $A = \sigma(C_{i})$, $\mathrm{tp}(a/A)=\sigma(r_{i})$, and $B \ind_{A} a$, equivalently $a \ind_A B$ by symmetry.  But then $\mathrm{SU}(a/A)= \infty$, so $\mathrm{SU}(a/AB) = \infty$, so $\mathrm{SU}(a/B)=\infty$, a contradiction.  By compactness, to show that the above set is consistent, it suffices to show, for any $B_{0} \subset B$ finite, that the part just mentioning $B_{0}$ and $a$ is consistent, or equivalently, that the part of $\Phi_{i}(B, Y ,a)$ mentioning just $B_{0}$ and $a$ is consistent for some $1 \leq i \leq m$.  But $\mathrm{tp}(a/B_{0})$ is $r_{i}$-resolvable for some $1 \leq i \leq m$, so there is some automorphism $\sigma$ of $\mathbb{M}$ so that, for $A = \sigma(C_{i})$, $\mathrm{tp}(a/A)=\sigma(r_{i})$, and $a \ind_{A} B_{0}$, equivalently $B_{0} \ind_A a$ by symmetry; as above, this is what consistency of that part of $\Phi_{i}(B, Y ,a)$ means.

\end{proof}

\end{document}
